\documentclass[11pt]{article}
\usepackage[a4paper,margin=0.92in]{geometry}
\usepackage{amsmath,amssymb,amsthm,mathtools}
\usepackage{graphicx}
\usepackage{microtype}
\usepackage[hidelinks]{hyperref}
\usepackage{enumitem}
\setlength{\parindent}{0pt}
\setlength{\parskip}{0.55em}

\newtheorem{theorem}{Theorem}
\newtheorem{lemma}{Lemma}
\newcommand{\Y}{\mathcal{Y}}
\newcommand{\TV}{\operatorname{TV}}
\newcommand{\GL}{\operatorname{GL}}

\title{No Extra Epsilon Scaling for Unbounded \(L^p\) Graphons\\
\large An \(L^1\) sharp-interface theorem and coefficient trichotomy}
\author{fa\_banach\_001, lane 4}
\date{13 August 2026}

\begin{document}
\maketitle

\textbf{Status.} Candidate full solution, likely valid; specialist review requested.

\section{Source question}

Zhang, Scott, Du, and Porter \cite{ZSDP} define a graphon Ginzburg--Landau functional on Young measures and prove its sharp-interface limit for bounded graphons. Section~6.2 of arXiv:2408.00422v2 asks how the Dirichlet term should be scaled for unbounded \(L^p\) graphons, suggesting that the rate should depend on the singularity of \(W\). Their conclusion again states that this scaling was not determined for general \(L^p((0,1)^2)\) kernels. The source takes \(p\geq1\).

\begin{figure}[ht]
\centering
\includegraphics[width=0.98\textwidth]{figures/open_problem_crop.png}
\caption{The scaling question and the source's \(L^\infty\)-only theorem on page 18 of arXiv:2408.00422v2.}
\end{figure}

\section{Idea of the proof}

The apparent obstruction disappears once one distinguishes boundedness of the functional from finiteness on its limiting domain. Every \(L^p\) graphon in the source regime belongs to \(L^1\). On a binary Young measure, the phase difference is at most \(2\), so the Dirichlet energy is at most \(4\lVert W\rVert_1\), even when \(W\) is unbounded.

Global continuity of the Dirichlet energy under narrow convergence is neither true nor needed. Truncate the unbounded integrand \(|\lambda-\mu|^2\). Narrow convergence of product Young measures gives weak-star convergence for every bounded truncation, which can be tested against \(W\in L^1\). Letting the truncation level increase proves lower semicontinuity. The same device handles the quartic double-well potential. The liminf and constant recovery sequence then give the full Gamma-limit. This also classifies every convergent prefactor \(a_\epsilon\).

\section{Definitions and result}

Let \(X=(0,1)\) with Lebesgue measure. Let \(W:X^2\to[0,\infty]\) be measurable with \(W\in L^1(X^2)\). Write \(\Y\) for the source space of parametrized probability Young measures \(\nu=\{\nu_x\}_{x\in X}\), equipped with narrow convergence. Define the extended-valued energies
\begin{align*}
D^W(\nu)
&:=\int_{X^2}W(x,y)\int_{\mathbb R^2}|\lambda-\mu|^2
\,d\nu_x(\lambda)d\nu_y(\mu)\,dx\,dy,\\
V(\nu)
&:=\int_X\int_{\mathbb R}\Phi(\lambda)\,d\nu_x(\lambda)\,dx,
\end{align*}
where \(\Phi:\mathbb R\to[0,\infty)\) is continuous and
\[
\Phi^{-1}(0)=\{-1,1\}.
\]
This includes the source potential \(\Phi(t)=(t^2-1)^2\). Let
\[
\Y^b:=\{\nu\in\Y:\nu_x(\{-1,1\})=1\text{ for a.e. }x\}.
\]
For positive coefficients \(a_\epsilon\), set
\[
F_\epsilon(\nu):=a_\epsilon D^W(\nu)+\frac1\epsilon V(\nu).
\]

\begin{theorem}[Integrable-graphon scaling classification]\label{thm:main}
Let \(\epsilon\downarrow0\).
\begin{enumerate}[label=\textup{(\roman*)}]
\item If \(a_\epsilon\to a\in[0,\infty)\), then \(F_\epsilon\) Gamma-converges narrowly to
\[
F_a(\nu)=
\begin{cases}
aD^W(\nu),&\nu\in\Y^b,\\
+\infty,&\nu\notin\Y^b.
\end{cases}
\]
\item If \(a_\epsilon\to+\infty\), then \(F_\epsilon\) Gamma-converges narrowly to
\[
F_\infty(\nu)=
\begin{cases}
0,&\nu\in\Y^b\text{ and }D^W(\nu)=0,\\
+\infty,&\text{otherwise}.
\end{cases}
\]
\end{enumerate}
In particular, for \(a_\epsilon\equiv1\),
\[
\GL^W_\epsilon\ \mathop{\longrightarrow}^{\Gamma}\ \TV^W
\]
for every nonnegative \(W\in L^1(X^2)\), where \(\TV^W\) is exactly the source graphon total variation. Hence the source scaling \(1\) and \(1/\epsilon\) works for every \(L^p\) graphon with \(p\geq1\), without dependence on its singularity strength.
\end{theorem}

\section{Proof}

We first isolate the only continuity statement that is needed.

\begin{lemma}\label{lem:lsc}
Both \(D^W\) and \(V\) are lower semicontinuous for narrow convergence on \(\Y\).
\end{lemma}

\begin{proof}
Suppose \(\nu^n\to\nu\) narrowly. For \(R>0\), let
\[
q_R(\lambda,\mu):=\min\{|\lambda-\mu|^2,R\}.
\]
This is bounded and continuous. Narrow convergence of product Young measures \cite[Lemma~8]{BCD} implies that
\[
h_{n,R}(x,y):=\int_{\mathbb R^2}q_R(\lambda,\mu)
\,d\nu^n_x(\lambda)d\nu^n_y(\mu)
\]
converges weak-star in \(L^\infty(X^2)\) to the corresponding \(h_R\) for \(\nu\). Because \(W\in L^1(X^2)\),
\[
\int_{X^2}Wh_{n,R}\longrightarrow\int_{X^2}Wh_R.
\]
Moreover \(D^W(\nu^n)\geq\int Wh_{n,R}\). Therefore
\[
\liminf_{n\to\infty}D^W(\nu^n)\geq\int_{X^2}Wh_R.
\]
Letting \(R\uparrow\infty\) and using monotone convergence yields
\[
\liminf_{n\to\infty}D^W(\nu^n)\geq D^W(\nu).
\]
For \(V\), repeat the argument with the bounded continuous truncations
\(\Phi_R=\min\{\Phi,R\}\), now using the constant test function \(1\in L^1(X)\), and then let \(R\uparrow\infty\).
\end{proof}

We prove Theorem~\ref{thm:main}. Consider arbitrary sequences \(\epsilon_n\downarrow0\) and \(\nu^n\to\nu\) narrowly. Observe first that
\[
V(\nu)=0\quad\Longleftrightarrow\quad\nu\in\Y^b.
\]
Indeed, the integrand is nonnegative and its zero set is exactly \(\{-1,1\}\). If \(\nu\notin\Y^b\), then \(V(\nu)>0\). Lemma~\ref{lem:lsc} gives \(\liminf_nV(\nu^n)\geq V(\nu)\), so
\[
\liminf_{n\to\infty}\frac{V(\nu^n)}{\epsilon_n}=+\infty.
\]
This proves the Gamma-liminf inequality at every nonbinary \(\nu\), for all coefficient regimes.

Suppose now that \(\nu\in\Y^b\). If \(a_{\epsilon_n}\to a\in(0,\infty)\), then for each \(\delta\in(0,a)\), eventually \(a_{\epsilon_n}\geq a-\delta\). Lemma~\ref{lem:lsc} gives
\[
\liminf_nF_{\epsilon_n}(\nu^n)
\geq(a-\delta)D^W(\nu).
\]
Let \(\delta\downarrow0\). For \(a=0\), nonnegativity directly gives the required lower bound \(0\). If \(a_{\epsilon_n}\to\infty\) and \(D^W(\nu)>0\), lower semicontinuity implies that \(D^W(\nu^n)\geq D^W(\nu)/2\) for all sufficiently large \(n\); hence the liminf is infinite. If \(D^W(\nu)=0\), nonnegativity gives the required lower bound \(0\).

It remains to construct recovery sequences. If \(\nu\in\Y^b\), take the constant sequence \(\nu^n=\nu\). Then \(V(\nu)=0\), and
\[
D^W(\nu)\leq4\lVert W\rVert_{L^1(X^2)}<\infty.
\]
Thus \(F_{\epsilon_n}(\nu)=a_{\epsilon_n}D^W(\nu)\), which tends to the stated finite-coefficient limit; it is identically zero in the divergent-coefficient case whenever \(D^W(\nu)=0\). At points where the proposed limit is \(+\infty\), the Gamma-limsup requirement is automatic. This proves both parts.

Finally, for \(\nu\in\Y^b\), the only possible values of \(|\lambda-\mu|\) are \(0\) and \(2\), so
\[
|\lambda-\mu|^2=2|\lambda-\mu|.
\]
Consequently \(D^W(\nu)\) equals the source definition of \(\TV^W(\nu)\), proving the last assertion. \qed

\section{Scope and novelty check}

The theorem applies to nonnegative integrable kernels. On the finite-measure square, \(L^p\subset L^1\) for every \(p\geq1\), so it resolves the scaling question for the entire \(L^p\)-graphon class defined in the source. It does not address signed kernels or genuinely nonintegrable kernels such as \(|x-y|^{-1-2s}\); those require a different theory and are not \(L^p\) graphons in the stated regime.

The proof deliberately uses lower semicontinuity, not global continuity: moment escape prevents global continuity of \(D^W\) on all Young measures. Extended values away from \(\Y^b\) are harmless in Gamma convergence.

The bounded novelty search, completed 13 August 2026, covered the run's four lightweight indexes, arXiv:2408.00422v2 and its bibliography, exact-title citation searches, and arXiv-oriented searches combining graphon GL, \(L^1/L^p\), unbounded kernels, Gamma convergence, and epsilon scaling. No later resolution of this extension or coefficient trichotomy was located. The source's 2026 publication still states the problem.

\section{Human-review recommendation}

Priority specialist review is recommended. The core check is short: verify the product-narrow weak-star convergence for \(q_R\), the \(L^1\) testing step, and the two monotone-convergence passages. A reviewer should also confirm that the intended graphon convention is nonnegative; this is standard and is used throughout the source's energy arguments.

\begin{thebibliography}{9}
\bibitem{ZSDP}
E. J. Zhang, J. Scott, Q. Du, and M. A. Porter,
\emph{Ginzburg--Landau Functionals in the Large-Graph Limit},
arXiv:2408.00422v2 (2025); Pure and Applied Functional Analysis 11 (2026), 169--209.

\bibitem{BCD}
A. Braides, P. Cermelli, and S. Dovetta,
\emph{Gamma-limit of the cut functional on dense graph sequences},
ESAIM Control Optim. Calc. Var. 26 (2020), article 26.
\end{thebibliography}

\end{document}
